% !TEX program = xelatex
\documentclass[12pt, a4paper]{article}

% --- Пакеты для русского языка и кодировки ---
\usepackage{fontspec}
\usepackage{polyglossia}
\setdefaultlanguage{russian}
\setotherlanguage{english}
\setmainfont{Times New Roman}
\setsansfont{Arial}
\setmonofont{Courier New}

% --- Пакеты для полей и форматирования (ГОСТ) ---
\usepackage[left=30mm, right=15mm, top=20mm, bottom=20mm]{geometry}
\usepackage{indentfirst}
\usepackage{setspace}
\onehalfspacing
\usepackage{misccorr}

% --- Пакеты для таблиц, списков, ссылок ---
\usepackage{longtable}
\usepackage{multirow}
\usepackage{booktabs}
\usepackage{enumitem}
\setlist{leftmargin=25pt}
\usepackage{hyperref}
\hypersetup{
    colorlinks=true,
    linkcolor=black,
    urlcolor=black,
}
\usepackage{amsmath}
\usepackage{graphicx}
\graphicspath{{figures/}}

% --- Нумерация страниц ---
\usepackage{fancyhdr}
\pagestyle{fancy}
\fancyhf{}
\rfoot{\thepage}
\renewcommand{\headrulewidth}{0pt}

\begin{document}

% --- Титульный лист ---
\begin{titlepage}
    \begin{center}
        \vspace*{1cm}
        
        {\large \textbf{ИП Климук Д.И.}}
        
        \vspace{4cm}
        
        {\Large \textbf{УТВЕРЖДАЮ}}
        
        \vspace{0.5cm}
        
        {\large ИП \_\_\_\_\_\_\_\_\_\_\_\_\_\_\_\_\_\_ Климук Д.И.}
        
        \vspace{0.3cm}
        
        {\large «\_\_\_» \_\_\_\_\_\_\_\_\_\_ 2026 г.}
        
        \vspace{4cm}
        
        {\huge \textbf{ЭСКИЗНЫЙ ПРОЕКТ}}
        
        \vspace{0.5cm}
        
        {\Large на программу \enquote{Motion Visualizer}}
        
        \vspace{1cm}
        
        {\large Шифр темы: MV-PET-01-2026}
        
        \vspace{1cm}
        
        {\normalsize \textbf{Пояснительная записка}}
        
        \vfill
        
        \begin{tabular}{ll}
            Листов & \textbf{20} \\
        \end{tabular}
        
        \vspace{1cm}
        
        \begin{tabularx}{\textwidth}{Xr}
            \textbf{СОГЛАСОВАНО} & \textbf{РАЗРАБОТАНО} \\
            (при наличии заказчика) & \\
            & \\
            & \hfill \rule{5cm}{0.4pt} \\
            & \hfill \rule{5cm}{0.4pt} \\
            «\_\_\_» \_\_\_\_\_\_\_\_ 2026 г. & \hfill \rule{5cm}{0.4pt} \\
            & \hfill Д.И. Климук \\
            & «\_\_\_» \_\_\_\_\_\_\_\_ 2026 г.
        \end{tabularx}
    \end{center}
\end{titlepage}

% --- Содержание ---
\tableofcontents
\newpage

% --- Введение ---
\section{Введение}

\subsection{Наименование программы}
Полное наименование программы: \enquote{Программа для визуализации движения твердотельного объекта в трёхмерной среде с регистрацией параметров в базе данных}.

Краткое наименование: \enquote{Motion Visualizer}.

\subsection{Условное обозначение темы разработки}
MV-PET-01-2026.

\subsection{Основание для разработки}
Разработка выполнена в рамках самостоятельного задания (пет-проекта), утверждённого ИП Климук Д.И. \enquote{05} июня 2026 года.

\subsection{Цель разработки}
Создание демонстрационного прототипа, подтверждающего техническую возможность и эффективность связки:
\begin{itemize}
    \item графический интерфейс пользователя (PySide6);
    \item 3D-визуализация (OpenGL);
    \item математическое ядро на C++ (интегрированное через pybind11);
    \item база данных для хранения параметров и результатов (SQLite).
\end{itemize}

\section{Назначение и область применения}

\subsection{Назначение программы}
Программа предназначена для моделирования и визуализации движения объекта (материальной точки) по заданной траектории с регистрацией параметров движения в базе данных.

\subsection{Область применения}
Прототип может быть использован в качестве основы для следующих систем:
\begin{itemize}
    \item учебные тренажёры;
    \item симуляторы движения;
    \item цифровые полигоны для испытаний;
    \item системы мониторинга и анализа движения.
\end{itemize}

\section{Технические характеристики}

\subsection{Постановка задачи на разработку программы}

Требуется разработать программу, обеспечивающую выполнение следующих функций:
\begin{enumerate}
    \item запуск симуляции движения объекта из состояния покоя;
    \item движение объекта из начальной точки $A(0,0,0)$ в конечную точку $B(100,0,0)$ по линейной траектории;
    \item запись в базу данных времени старта и окончания движения;
    \item остановка объекта по достижении конечной точки;
    \item сброс объекта в начальную позицию по команде пользователя.
\end{enumerate}

\subsection{Описание применяемых математических методов}

Для моделирования движения используется метод интегрирования дифференциальных уравнений движения материальной точки.

Уравнение движения в векторной форме:
\begin{equation}
    m \frac{d\mathbf{v}}{dt} = \mathbf{F}_{\text{сум}}
\end{equation}

В скалярной форме для одномерного движения вдоль оси $x$:
\begin{equation}
    \frac{dx}{dt} = v(t), \quad \frac{dv}{dt} = a(t)
\end{equation}

В начальный момент времени $t_0$:
\begin{equation}
    x(t_0) = 0, \quad v(t_0) = 0
\end{equation}

При постоянном ускорении $a$:
\begin{equation}
    x(t) = x_0 + v_0 t + \frac{at^2}{2}
\end{equation}

Для численного интегрирования предлагается использовать метод Эйлера:
\begin{align}
    v_{n+1} &= v_n + a \cdot \Delta t \\
    x_{n+1} &= x_n + v_n \cdot \Delta t
\end{align}

\subsection{Допущения и ограничения}
\begin{enumerate}
    \item объект считается материальной точкой (размерами и формой пренебрегаем);
    \item движение происходит в одной плоскости без учёта вращения;
    \item внешние силы (трение, сопротивление среды) не учитываются;
    \item шаг интегрирования $\Delta t$ принимается равным 0,1 секунды.
\end{enumerate}

\subsection{Описание алгоритма решения задачи}

Общий алгоритм работы программы представлен на рисунке 1.

\begin{figure}[h]
    \centering
    \framebox(400,300)[c]{\parbox{12cm}{\centering \textbf{Место для блок-схемы алгоритма} \\ (будет добавлена на этапе технического проекта)}}
    \caption{Общая блок-схема алгоритма работы программы}
    \label{fig:algorithm}
\end{figure}

\subsection{Обоснование выбора метода организации входных и выходных данных}

\subsubsection{Входные данные}
Входными данными являются параметры, задаваемые пользователем через графический интерфейс:
\begin{itemize}
    \item начальные координаты (по умолчанию $x=0$);
    \item конечные координаты (по умолчанию $x=100$);
    \item скорость движения (по умолчанию $v=5$ ед./с).
\end{itemize}

\subsubsection{Выходные данные}
Выходные данные сохраняются в базе данных SQLite:
\begin{itemize}
    \item идентификатор сессии (первичный ключ);
    \item время начала движения;
    \item время окончания движения;
    \item пройденное расстояние;
    \item длительность движения.
\end{itemize}

\subsection{Описание и обоснование выбора состава технических и программных средств}

\subsubsection{Технические средства}
\begin{table}[h]
    \centering
    \begin{tabular}{|p{4cm}|p{3cm}|p{8cm}|}
        \hline
        \textbf{Компонент} & \textbf{Требование} & \textbf{Обоснование} \\
        \hline
        Процессор & x86\_64, 2+ ядра & Обеспечение производительности \\
        \hline
        Оперативная память & 4+ ГБ & Достаточно для работы приложения и СУБД \\
        \hline
        Видеокарта & Поддержка OpenGL 3.3+ & Требование Qt OpenGL \\
        \hline
        ОС & Ubuntu 22.04/24.04 & Стабильность, совместимость \\
        \hline
    \end{tabular}
    \caption{Требования к техническим средствам}
\end{table}

\subsubsection{Программные средства}
\begin{table}[h]
    \centering
    \begin{tabular}{|p{4cm}|p{3cm}|p{8cm}|}
        \hline
        \textbf{Компонент} & \textbf{Версия} & \textbf{Обоснование} \\
        \hline
        Язык Python & 3.13 & Современный, поддержка типизации \\
        \hline
        PySide6 & 6.8+ & Qt6 для Python, лицензия LGPL \\
        \hline
        PyOpenGL & 3.1+ & Стандартная обёртка OpenGL \\
        \hline
        Язык C++ & 23 & Высокая производительность \\
        \hline
        pybind11 & 2.13+ & Связь C++ и Python \\
        \hline
        SQLite & 3 & Встроенная БД, не требует сервера \\
        \hline
        VS Code & — & Удобная среда разработки \\
        \hline
    \end{tabular}
    \caption{Выбор программных средств}
\end{table}

\section{Ожидаемые технико-экономические показатели}

\subsection{Трудоёмкость разработки}
Планируемая трудоёмкость разработки составляет 10–12 человеко-часов.

\subsection{Эффективность}
Создание прототипа позволит сократить время на техническую демонстрацию компетенций при переговорах с потенциальными заказчиками. Ожидаемый экономический эффект — возможность заключения контрактов на сумму от 2,4 млн рублей.

\section{Источники, использованные при разработке}

\begin{enumerate}
    \item ГОСТ 19.102-77. Единая система программной документации. Стадии разработки.
    \item ГОСТ 19.404-79. Единая система программной документации. Пояснительная записка. Требования к содержанию и оформлению.
    \item Qt Documentation. Qt for Python. \url{https://doc.qt.io/qtforpython/}
    \item pybind11 Documentation. \url{https://pybind11.readthedocs.io/}
    \item Learn OpenGL. \url{https://learnopengl.com/}
\end{enumerate}

\end{document}